\section{Method}

\subsection{Overview}

DreamerVLA couples a one-trajectory OpenVLA-OFT policy with a token-space
world model and a success classifier.  The system uses one external observation
representation throughout collection, replay, warm-up, and online cotraining:
the projected visual input tokens produced before the language-model action
positions.  For each environment step $t$,
\[
E_t = f_{\mathrm{OFT}}^{\mathrm{vision}}(I_t,\ell)
    \in \mathbb{R}^{256\times4096},
\]
where $I_t$ is the current agent-view image and $\ell$ is the task instruction.
The 256 visual tokens are persisted under the HDF5 key
\texttt{obs\_embedding}; action-decoder states are never used as world-model
observations.

The trainable system contains three components: (i) a chunk-aware Transformer
world model that predicts future input-token grids, (ii) a spatial success
classifier trained on windows of the same token grids, and (iii) an actor bridge
that maps predicted input tokens to OpenVLA-OFT's internal discrete action
positions.  The bridge creates action positions only inside the decoder; they
are not stored in replay and do not define a second observation interface.

\subsection{Cold-Start and Online Cotraining Pipeline}

Training follows four stages.  First, the one-trajectory OpenVLA-OFT checkpoint
collects real LIBERO rollouts.  Each rollout produces a reward shard and a
shape-aligned input-token sidecar.  Before reuse, the loader verifies the source
name, token count, token width, storage format, and every per-demo tensor shape.

Second, collected episodes seed a shared online replay.  Third, the world model
and success classifier are warmed up from this replay using exactly the update
functions used online.  Finally, real and imagined trajectories are interleaved
during cotraining.  In the synchronous route one runner owns the loop.  In the
asynchronous route, the actor group owns FSDP policy updates, the rollout group
owns no-gradient policy inference, the environment group owns real and imagined
stepping, and the learner group owns world-model and classifier updates.

\subsection{Input-Token Observation Contract}

For an episode of length $T$, the sidecar stores
\[
\mathcal{E}=[E_1,\ldots,E_T]
    \in \mathbb{R}^{T\times256\times4096}.
\]
The contract fixes one image, one-frame policy history, 256 patches, and a token
width of 4096.  Its flattened size, $1{,}048{,}576$, is metadata only; model
interfaces retain the two token axes.  Reward and sidecar shards must have the
same files, demonstrations, and frame lengths.  This validation is performed
before collection resume, replay seeding, standalone classifier training, and
standalone world-model training.

Language and proprioception remain distinct conditioning signals.  They are
projected to compact embeddings and concatenated to token channels inside the
world model and classifier; they do not alter the persisted visual-token
schema.

\subsection{Chunk-Aware Token World Model}

Let $a_t\in\mathbb{R}^{7}$ denote the robot action and let $K=8$ be the action
chunk length.  The world model embeds the action and concatenates it to every
visual token.  Optional language and proprioceptive embeddings are concatenated
in the same role-separated manner.  Given a history of $H$ token grids, a
Transformer transition predicts the next grid,
\[
\widehat{E}_{t+1}
  = F_\theta(E_{t-H+1:t},a_t,p_t,\ell).
\]
Multi-step prediction is autoregressive: each predicted grid is appended to
the history before the next action is applied.  For $J$ predicted chunks, the
training loss is
\[
\mathcal{L}_{\mathrm{WM}}
 = \sum_{j=1}^{JK} \alpha_j
   \left\|\widehat{E}_{t+j}-E_{t+j}\right\|_2^2
   + \lambda_{\cos}
   \bigl(1-\cos(\widehat{E}_{t+j},E_{t+j})\bigr),
\]
with optional auxiliary reward and proprioceptive reconstruction terms.  The
mainline keeps the visual tokens explicit rather than accepting flattened
observation vectors.

\subsection{Spatial Success Classifier}

The classifier receives a window
$E_{t-W+1:t}\in\mathbb{R}^{W\times256\times4096}$.
Each visual token is projected to the classifier width, augmented with temporal
and spatial position embeddings, and processed by a Transformer encoder with a
classification token.  A binary head estimates
\[
q_\phi(y_t=1\mid E_{t-W+1:t},p_{t-W+1:t},\ell),
\]
where $y_t$ indicates task success.  Warm-up and online updates share the same
balanced WMPO sampling protocol and binary cross-entropy objective.  During
imagined rollout, this score supplies the success signal used by policy
optimization.

\subsection{OpenVLA Discrete-Action Bridge}

The actor consumes a predicted token grid
$\widehat{E}_t\in\mathbb{R}^{256\times4096}$.  A Transformer decoder attends a
set of $K\times7=56$ learned action positions to the 256 source tokens, producing
4096-dimensional states only inside the OpenVLA decoder.  The frozen or
finetuned language-model head converts those states into categorical action
tokens, which are mapped to continuous controls.  For action dimension $d$ and
chunk step $k$,
\[
\pi_\psi(z_{k,d}\mid\widehat{E}_t)
  = \operatorname{softmax}\bigl(W_{\mathrm{LM}}h_{k,d}\bigr).
\]
The policy objective operates on token-level log probabilities.  This decoder
construction preserves OpenVLA-OFT's action vocabulary while maintaining a
single, unambiguous world-model observation state.

\subsection{Cotraining Objective}

World-model and classifier parameters are optimized from replay, while actor
parameters are optimized from grouped real and imagined trajectories.  The
actor update combines relative advantages, clipped likelihood ratios, entropy
regularization, and a reference-policy KL penalty.  Parameter ownership is
explicit: the rollout replica periodically pulls actor weights, and imagined
environment workers periodically receive world-model and classifier weights.
All reported task quality is measured in the real LIBERO environment; imagined
success is a training signal rather than an evaluation substitute.
